% Provenance (never printed): sources/book2/BRIEFS.md, g11-11. Lorenz curve and Gini, interdecile ratio and top shares,
% before and after taxes and transfers, income and wealth, skills and technical change, mobility, what the measures miss.
% Real numbers: figdata/grade-11/11-inequality/ (Eurostat ilc_di12 and ilc_di12b, 2025; PIP percentiles; Brenland's
% household survey, invented). Ledger: Piketty and Saez 2003 (Table II), Autor 2014, Chetty et al. 2014, Eurostat
% equivalence scale. The third chart (Gini against income across countries) was not drawn (page budget).
\chapter{Inequality}\label{ch:g11:inequality}

In 2025, before pensions and other social transfers, the \omterm{def:g11:inequality:gini}{Gini coefficient} of income in France was 53.7 on a scale from 0 to 100; after them it was 30.4. In
Slovakia the figures were 39.9 and 23.0. Every European country in the chart below is far less unequal after its transfers than before, in the accounts; the gap ranges from 13.8 points
in Spain to 23.3 in France.
% figdata: eu_gini (FR 53.7 -> 30.4, cut 23.3; SK 39.9 -> 23.0; ES cut 13.8)

How unequal a country is depends on what is measured, among whom, and before or after the state's taxes and transfers. This chapter builds the standard measures, the
\omterm{def:g11:inequality:lorenz}{Lorenz curve} and the \omterm{def:g11:inequality:gini}{Gini coefficient} above all, shows how they are computed and what they leave out, and looks at two questions the evidence can answer: why the pay
of the educated pulled away from the rest, and how far children's incomes follow their parents'.

\begin{openingchart}{Gini coefficient of equivalised disposable income, before social transfers (pensions counted as transfers) and after, twelve European countries and
the European Union, survey year 2025. Scale 0 (all equal) to 100.}{Eurostat, EU statistics on income and living conditions (ilc\_di12, ilc\_di12b)}
\begin{tikzpicture}
\begin{axis}[width=0.92\linewidth, height=5.8cm, ybar, /pgf/bar width=5pt, xtick={0,...,12}, xticklabels={SK, BE, SI, PL, NL, SE, DE, FR, IT, ES, LT, BG, EU},
  enlarge x limits=0.05, ymin=0, ymax=62, ylabel={Gini (0--100)}, legend style={at={(0.02,0.98)}, anchor=north west, font=\scriptsize, draw=none},
  legend columns=2, x tick label style={/tikz/font=\scriptsize}]
\addplot[fill=omThm!60, draw=omThm] table[x=x, y=before] {figdata/grade-11/11-inequality/out/eu_gini.csv};
\addplot[fill=omDef!60, draw=omDef] table[x=x, y=after] {figdata/grade-11/11-inequality/out/eu_gini.csv};
\legend{before social transfers, after}
\end{axis}
\end{tikzpicture}
\end{openingchart}

\begin{recall}
From the grades 6--9 volume: mean and median income, the shares of income received by each fifth of households, disposable income, the world distribution of income,
and the poverty line. From the mathematics volumes: quartiles and percentiles. From \cref{ch:g11:marginal-reasoning}: marginal and average tax rates.
\end{recall}

\section{The Lorenz curve and the Gini coefficient}\label{sec:g11:inequality:lorenz}

\begin{definition}[Lorenz curve]\label{def:g11:inequality:lorenz}
Rank households from the poorest to the richest. The \emph{Lorenz curve}\index{Lorenz curve} gives, for each share $x$ of households counted from the bottom, the
share $L(x)$ of total income they receive.
\end{definition}

\begin{definition}[Gini coefficient]\label{def:g11:inequality:gini}
The \emph{Gini coefficient}\index{Gini coefficient} is the area between the diagonal and the \omterm{def:g11:inequality:lorenz}{Lorenz curve}, divided by the whole area under the diagonal:
$G = 1 - 2B$, where $B$ is the area under the \omterm{def:g11:inequality:lorenz}{Lorenz curve}.
\end{definition}

\begin{proposition}[Properties of the Lorenz curve and the Gini]\label{prop:g11:inequality:properties}
The \omterm{def:g11:inequality:lorenz}{Lorenz curve} runs from $(0,0)$ to $(1,1)$, is convex, and lies on or below the diagonal; $0 \le G < 1$, and $G = 0$ exactly when all incomes are equal.
\end{proposition}
\begin{proof}
The slope of the curve between two points is the income of the households added, divided by their number, relative to the mean: their average income over the mean
income. Since households are ranked from the poorest, this slope never falls as $x$ rises, so the curve is convex. A convex curve from $(0,0)$ to $(1,1)$ lies on or below
the straight line joining them, so $B \le \tfrac12$ and $G \ge 0$. If all incomes are equal, every slope is 1, the curve is the diagonal, $B = \tfrac12$ and $G = 0$.
Otherwise the poorest have a slope below 1 at the start and the curve dips below the diagonal, so $G > 0$. Since incomes are not negative, $L \ge 0$, and $B > 0$ when
some income is positive, so $G < 1$.
\end{proof}

\begin{measure}[How a Gini coefficient is computed from a survey]
With incomes grouped into fifths, the \omterm{def:g11:inequality:lorenz}{Lorenz curve} is known at $x = 0, 0.2, \dots, 1$, with cumulative shares $L_0 = 0, L_1, \dots, L_5 = 1$. Join the points by straight
lines; the area under them is a sum of trapezoids, $B = \sum_{k=1}^{5} 0.2 \times (L_{k-1} + L_k)/2$, and
\[ G = 1 - \sum_{k=1}^{5} 0.2 \times (L_{k-1} + L_k). \]
Straight lines lie above the true curve between the points, since the curve is convex, so grouping hides the inequality inside each group and gives a Gini slightly
too low. Statistical offices compute it from each household's income.
\end{measure}

\begin{omfigure}
\begin{tikzpicture}
\begin{axis}[width=0.46\linewidth, height=0.46\linewidth, xmin=0, xmax=100, ymin=0, ymax=100, xtick={0,20,...,100}, ytick={0,20,...,100},
  xlabel={households, poorest first (\%)}, ylabel={share of income (\%)}, grid=major, grid style={black!10}, tick label style={/tikz/font=\scriptsize},
  legend style={at={(0.03,0.97)}, anchor=north west, font=\scriptsize, draw=none, fill=none}, legend cell align=left]
\addplot[name path=diag, black, thin, forget plot] coordinates {(0,0) (100,100)};
\addplot[name path=mkt, omThm, thick, mark=*, mark size=1.2pt] table[x=pop, y=market] {figdata/grade-11/11-inequality/out/bren_lorenz.csv};
\addplot[omDef, thick, mark=*, mark size=1.2pt] table[x=pop, y=disposable] {figdata/grade-11/11-inequality/out/bren_lorenz.csv};
\addplot[omThm!20, forget plot] fill between[of=diag and mkt];
\legend{market income, disposable income}
\end{axis}
\end{tikzpicture}

{\small \omterm{def:g11:inequality:lorenz}{Lorenz curves} of Brenland's household survey (invented), income divided by the square root of household size, by fifths. The shaded area between the diagonal
and the market-income curve, divided by $\tfrac12$, is its \omterm{def:g11:inequality:gini}{Gini coefficient}, 0.30; for disposable income it is 0.24.}
\end{omfigure}

\begin{definition}[Interdecile ratio and top income share]\label{def:g11:inequality:interdecile}
The \emph{interdecile ratio}\index{interdecile ratio} is the income at the 90th percentile divided by the income at the 10th. A \emph{top income share}\index{top income
share}\phantomsection\label{def:g11:inequality:top-share} is the share of total income received by the richest group, such as the top tenth or the top 1\,\%.
\end{definition}

The Gini summarises the whole curve in one number, and different curves can give the same Gini. The \omterm{def:g11:inequality:interdecile}{interdecile ratio} ignores the very top and the very bottom; top
shares look only at the top, where household surveys capture incomes badly and tax records do better.

\begin{record}[Top income shares in the twentieth century]
% ledger: g11:ps2003-top1, g11:ps2003-top10, g11:autor2014-top1
Tax returns of the United States show the share of income, excluding capital gains, received by the top 1\,\% of tax units: 18.0\,\% in 1913, 19.6\,\% in 1928, 10.5\,\%
in 1944, 7.8\,\% in 1970 and 14.6\,\% in 1998. The share of the top tenth stood at 40 to 45\,\% between the wars, fell to about 30\,\% during the Second World War, stayed at
31 to 32\,\% until the 1970s, and was above 40\,\% again by the mid-1990s. A household measure of the top 1\,\% share went from 10.0\,\% in 1979 to 22.5\,\% in 2012. The
pattern is a U: high, then low for three decades, then high again.
\end{record}

\section{Before and after taxes and transfers}\label{sec:g11:inequality:redistribution}

\begin{definition}[Market income]\label{def:g11:inequality:market-income}
A household's \emph{market income}\index{market income} is its income from work, property and business, before taxes and before transfers from the state.
\end{definition}

\begin{definition}[Progressive and regressive taxes]\label{def:g11:inequality:progressive}
A tax is a \emph{progressive tax}\index{progressive tax} when its average rate rises with income, and a \emph{regressive tax}\index{regressive
tax}\phantomsection\label{def:g11:inequality:regressive} when its average rate falls as income rises.
\end{definition}

\begin{definition}[Equivalence scale]\label{def:g11:inequality:equivalence-scale}
An \emph{equivalence scale}\index{equivalence scale} is the number by which a household's income is divided to compare households of different sizes; it is smaller than
the number of members because people living together share costs.
\end{definition}

\begin{measure}[Equivalence scales]
% ledger: g11:eurostat-equiv
Eurostat divides a household's income by its size on the modified OECD scale: 1 for the first adult, 0.5 for each other person aged 14 or over, 0.3 for each child under 14.
A couple with two young children counts as $1 + 0.5 + 0.3 + 0.3 = 2.1$. Brenland's survey uses the square root of the number of members, 2 for a household of four. The
choice matters: a scale that gives children less weight makes large families look richer.
\end{measure}

The measured effect of taxes and transfers depends on what is counted. Eurostat's \emph{before} figure in the opening chart removes pensions and other social transfers
but not income taxes; in Brenland's survey, \emph{market} income is before both. Pensions do much of the work in the European figures: a retired household has almost no
\omterm{def:g11:inequality:market-income}{market income} and a pension that replaces it. The \emph{before} figure is an accounting one: it is not the inequality the country
would have without its transfers, since people would then save and work differently, and retired households would not live on nothing.

\section{Income, wealth and skills}\label{sec:g11:inequality:wealth}

\begin{definition}[Wealth]\label{def:g11:inequality:wealth}
A household's \emph{wealth}\index{wealth} is the value of what it owns (a house, savings, shares, a business) minus its debts, at a date.
\end{definition}

Income is a flow over a year; \omterm{def:g11:inequality:wealth}{wealth} is a stock at a date. \omterm{def:g11:inequality:wealth}{Wealth} can be far more unequal than income, because it accumulates over a lifetime and is inherited, and
because many households own little more than they owe.

\begin{omchart}{Lorenz curves of income or consumption per person, latest survey year, four countries: Sweden (2023), the United States (2024), Brazil (2024), South Africa
(2022, consumption). The Gini coefficients from these percentiles are 0.29, 0.42, 0.50 and 0.54.}{World Bank, Poverty and Inequality Platform, percentiles (2021 PPP)}
\begin{tikzpicture}
\begin{axis}[width=0.5\linewidth, height=0.5\linewidth, xmin=0, xmax=100, ymin=0, ymax=100, xtick={0,20,...,100}, ytick={0,20,...,100},
  xlabel={people, poorest first (\%)}, ylabel={share (\%)}, grid=major, grid style={black!10}, tick label style={/tikz/font=\scriptsize},
  legend style={at={(0.03,0.97)}, anchor=north west, font=\scriptsize, draw=none, fill=none}, legend cell align=left]
\addplot[black, thin, domain=0:100, samples=2, forget plot] {x};
\addplot[omDef, thick] table[x=pop, y=cum] {figdata/grade-11/11-inequality/out/lorenz_SWE.csv};
\addplot[omProp, thick] table[x=pop, y=cum] {figdata/grade-11/11-inequality/out/lorenz_USA.csv};
\addplot[omMeth, thick] table[x=pop, y=cum] {figdata/grade-11/11-inequality/out/lorenz_BRA.csv};
\addplot[omThm, thick] table[x=pop, y=cum] {figdata/grade-11/11-inequality/out/lorenz_ZAF.csv};
\legend{Sweden, United States, Brazil, South Africa}
\end{axis}
\end{tikzpicture}
\end{omchart}
% figdata: lorenz_gini (SWE 2023 0.29, USA 2024 0.42, BRA 2024 0.50, ZAF 2022 consumption 0.54)

\begin{model}{Supply and demand for skills}
\assumes Two kinds of labour, with and without a college degree; technical change raises the demand for the first relative to the second; the supply of graduates grows
  as people study longer.
\predicts The graduates' wage premium rises when demand grows faster than supply, and stops rising when supply catches up.
\tested Consistent with it: in the United States the average college graduate earned 48\,\% more a week than the average high-school graduate in 1981, when the growth of the supply of
  graduates had slowed; 72\,\% more in 1990, 90\,\% in 2000 and 97\,\% in 2005. When the supply of new graduates rose fast between 2004 and 2012, the premium stopped
  rising (Autor 2014).
\breaks It explains the gap between graduates and the rest, not the rise of the top 1\,\% share, which concerns a few very high incomes rather than the two broad
  groups the model compares.
\end{model}
% ledger: g11:autor2014-premium

\section{From one generation to the next}\label{sec:g11:inequality:mobility}

\begin{definition}[Intergenerational mobility]\label{def:g11:inequality:mobility}
\emph{Intergenerational mobility}\index{intergenerational mobility} is the degree to which children's positions in the income distribution differ from their parents'.
\end{definition}

\begin{definition}[Mobility table]\label{def:g11:inequality:mobility-table}
A \emph{mobility table}\index{mobility table} counts families by the parents' position (rows) and the child's position (columns). Divided by the row totals, it gives
\emph{outflow} shares: where children of a given origin end up. Divided by the column totals, it gives \emph{inflow} shares: where people in a given position came from.
\end{definition}

\begin{omfigure}
\renewcommand{\arraystretch}{1.25}
\begin{tabular}{lrrrr}
\toprule
Parents' third $\backslash$ child's third & bottom & middle & top & total \\
\midrule
bottom & 50 & 30 & 20 & 100 \\
middle & 30 & 40 & 30 & 100 \\
top & 20 & 30 & 50 & 100 \\
\midrule
total & 100 & 100 & 100 & 300 \\
\bottomrule
\end{tabular}

{\small An invented \omterm{def:g11:inequality:mobility-table}{mobility table} of 300 families. Outflow: of the 100 children of bottom-third parents, $20/100 = 20\,\%$ reach the top third. Inflow: of the 100
people in the top third, 20\,\% have bottom-third parents and 50\,\% top-third parents.}
\end{omfigure}

Large studies use tax records. In the United States, records of more than 40 million children and their parents show that a rise of 10 percentiles in the parents' income
goes with a rise of about 3.4 percentiles in the child's, on average. The chance that a child born in the bottom fifth reaches the top fifth varied from 4.4\,\% in
Charlotte to 12.9\,\% in San Jose. These are correlations: they say how much position is passed on, not by what means.
% ledger: g11:chks2014

\section{What the measures miss}\label{sec:g11:inequality:limits}

Income measures miss what governments provide in kind, such as schooling and health care, and taxes on spending that surveys do not record. A household's income in one
year is a snapshot: students and the retired have low incomes but may be well off over their lives, so inequality of consumption, which people smooth by saving and
borrowing, can be lower than inequality of yearly income. Surveys miss the richest, whom tax records catch. And a household measure says nothing of how income is shared
inside the household.

\section{Exercises}

\begin{exercise}[$\star$]\label{exo:g11:inequality:1}
The fifths of a country's households receive 5, 10, 15, 25 and 45\,\% of income, poorest first. Give the points of its \omterm{def:g11:inequality:lorenz}{Lorenz curve}.
\end{exercise}

\begin{exercise}[$\star$]\label{exo:g11:inequality:2}
Same question for shares of 8, 12, 18, 24 and 38\,\%.
\end{exercise}

\begin{exercise}[$\star$]\label{exo:g11:inequality:3}
From the opening chart, which country has the highest Gini after transfers, and which has the largest fall from before to after? What are the figures for the EU as a
whole?
\end{exercise}

\begin{exercise}[$\star$]\label{exo:g11:inequality:4}
The 10th percentile of income is 15\,000 dalens and the 90th is 60\,000. Compute the \omterm{def:g11:inequality:interdecile}{interdecile ratio}.
\end{exercise}

\begin{exercise}[$\star$]\label{exo:g11:inequality:5}
In the invented \omterm{def:g11:inequality:mobility-table}{mobility table}, what share of the children of top-third parents end in the bottom third? What share of the bottom third have bottom-third parents?
\end{exercise}

\begin{exercise}[$\star\star$]\label{exo:g11:inequality:6}
Compute the \omterm{def:g11:inequality:gini}{Gini coefficient} of \cref{exo:g11:inequality:1} by trapezoids.
\end{exercise}

\begin{exercise}[$\star\star$]\label{exo:g11:inequality:7}
Compute the \omterm{def:g11:inequality:gini}{Gini coefficient} of \cref{exo:g11:inequality:2} by trapezoids. Which country is more equal?
\end{exercise}

\begin{exercise}[$\star\star$]\label{exo:g11:inequality:8}
Three Brenland households have \omterm{def:g11:inequality:market-income}{market incomes} of 8000, 30\,000 and 90\,000 dalens. They pay taxes of 0, 3000 and 21\,000 and receive transfers of 6000, 1000 and 0.
Compute their disposable incomes, and the ratio of the richest to the poorest before and after.
\end{exercise}

\begin{exercise}[$\star\star$]\label{exo:g11:inequality:9}
Households with incomes of 10\,000, 30\,000 and 100\,000 dalens pay an income tax of 500, 3000 and 15\,000, and a tax on purchases of 900, 1800 and 4000. Is each tax
progressive or regressive?
\end{exercise}

\begin{exercise}[$\star\star$]\label{exo:g11:inequality:10}
A \omterm{def:g11:inequality:mobility-table}{mobility table} of 300 families has rows (parents bottom, middle, top) of 50, 30, 20; 30, 45, 25; 20, 25, 55 (children bottom, middle, top). Compute the outflow shares of
the bottom row and the inflow shares of the top column.
\end{exercise}

\begin{exercise}[$\star\star$]\label{exo:g11:inequality:11}
Five households have incomes of 1000, 2000, 3000, 4000 and 10\,000. Compute the Gini by trapezoids (each household is a fifth). Then every household receives 1000 more:
compute it again.
\end{exercise}

\begin{exercise}[$\star\star\star$]\label{exo:g11:inequality:12}
What evidence would contradict the supply-and-demand-for-skills model? Is the halt of the college premium after 2004 such evidence? Is the rise of the top 1\,\% share?
\end{exercise}

\begin{exercise}[$\star\star\star$]\label{exo:g11:inequality:13}
Five households have incomes of 10, 20, 30, 40 and 100. After a reform they have 9, 25, 35, 41 and 90. Show that the Gini falls although the poorest household loses.
\end{exercise}

\begin{exercise}[$\star\star\star$]\label{exo:g11:inequality:14}
Household A is one adult with 20\,000 dalens; household B is a couple with two children under 14 and 36\,000. Rank them by total income, by income per person, by income on
the square-root scale and on the modified OECD scale.
\end{exercise}

\begin{exercise}[$\star\star\star$]\label{exo:g11:inequality:15}
Explain why the inequality of consumption is usually lower than the inequality of yearly income, and why South Africa's curve in the chart of four countries is not
strictly comparable with the others.
\end{exercise}

\section{Problem: Brenland's Household Survey}

\begin{problem}[{Weekend problem --- how much do taxes and transfers reduce inequality?}]
\label{pb:g11:inequality:1}
Brenland's survey of 2000 households records each household's \omterm{def:g11:inequality:market-income}{market income} and its disposable income. Divided by the square root of household size and ranked, the
fifths of households receive these shares (\%):

\medskip
\begin{center}
\begin{tabular}{lrrrrr}
\toprule
fifth & 1 & 2 & 3 & 4 & 5 \\
\midrule
\omterm{def:g11:inequality:market-income}{market income} & 6.9 & 12.8 & 17.5 & 23.6 & 39.2 \\
disposable income & 10.0 & 13.9 & 17.7 & 22.9 & 35.5 \\
\bottomrule
\end{tabular}
\end{center}
The 10th and 90th percentiles are 7800 and 33\,900 dalens for \omterm{def:g11:inequality:market-income}{market income} and 9300 and 28\,600 for disposable income; the top tenth receives 23.6\,\% of \omterm{def:g11:inequality:market-income}{market income} and
20.8\,\% of disposable income.
% figdata: bren_fifths, bren_key

\medskip
\noindent\textbf{Part I --- Comparing households.}
\begin{enumerate}
\item A household of four has 36\,000 dalens. What is its income on the square-root scale? Compare with a single person with 18\,000.
\item Why divide by the square root of the size rather than by the size itself?
\item What is the four-person household's income on the modified OECD scale, if two of them are children under 14?
\end{enumerate}

\medskip
\noindent\textbf{Part II --- Market income.}
\begin{enumerate}[resume]
\item Give the points of the \omterm{def:g11:inequality:lorenz}{Lorenz curve} of \omterm{def:g11:inequality:market-income}{market income}.
\item Compute the area under it by trapezoids.
\item Compute the \omterm{def:g11:inequality:gini}{Gini coefficient} of \omterm{def:g11:inequality:market-income}{market income}.
\item How many times the bottom fifth's share is the top fifth's?
\end{enumerate}

\medskip
\noindent\textbf{Part III --- Disposable income.}
\begin{enumerate}[resume]
\item Give the points of the \omterm{def:g11:inequality:lorenz}{Lorenz curve} of disposable income.
\item Compute its \omterm{def:g11:inequality:gini}{Gini coefficient}.
\item By how many points do taxes and transfers lower the Gini?
\item Which fifths gain share, and which lose?
\item Compare the fall with the falls in the opening chart, and say why the two are not measured the same way.
\end{enumerate}

\medskip
\noindent\textbf{Part IV --- Other measures.}
\begin{enumerate}[resume]
\item Compute the \omterm{def:g11:inequality:interdecile}{interdecile ratio} before and after.
\item By how much does the top tenth's share fall?
\item Compute the fall of each measure (Gini, \omterm{def:g11:inequality:interdecile}{interdecile ratio}, top tenth's share) in per cent of its value before. Which falls most?
\item Why is a Gini computed over fifths lower than one computed from each household's income?
\item Name two things the survey's disposable income leaves out.
\item A flat payment of the same amount to every household, paid for by a tax at the same rate on every income, is added. Does the Gini of disposable income rise or fall?
\item Can the Gini fall while the poorest households lose? Refer to the exercises.
\item State the result of the problem: the \omterm{def:g11:inequality:gini}{Gini coefficient} before and after taxes and transfers, and the reduction in points.
\end{enumerate}
\end{problem}
